\documentclass[a4,11pt]{aleph-notas}
% Se puede ver la documentación aquí:
% https://github.com/alephsub0/LaTeX_aleph-notas

% -- Paquetes adicionales
\usepackage{enumitem}
\usepackage{url}
\usepackage{array}
\usepackage{booktabs}
\usepackage{longtable}
\hypersetup{
    urlcolor=blue,
    linkcolor=blue,
}

% -- Datos
\institucion{Facultad de Ciencias Exactas, Naturales y Ambientales}
\carrera{Catálogo STEM}
\asignatura{Álgebra lineal}
\tema{Clase 03: Matriz inversa}
\autor{Andrés Merino}
\fecha{Periodo 2026-01}

\logouno[0.16\textwidth]{Logos/logoPUCE_04_ac}
\definecolor{colortext}{HTML}{1957d1}
\definecolor{colordef}{HTML}{1957d1}
\fuente{montserrat}

\begin{document}

\encabezado

%%%%%%%%%%%%%%%%%%%%%%%%%%%%%%%%%%%%%%%%
\section*{Resultado de Aprendizaje}
%%%%%%%%%%%%%%%%%%%%%%%%%%%%%%%%%%%%%%%%

%%%%%%%%%%%%%%%%%%%%%%%%%%%%%%%%%%%%%%%%
\subsection*{RdA de la asignatura:}
\begin{itemize}[leftmargin=*]
    \item \textbf{RdA 1:} Comprender los conceptos fundamentales del Álgebra Lineal, incluyendo el estudio de matrices, determinantes, sistemas de ecuaciones lineales y espacios vectoriales, destacando su importancia en el análisis y resolución de problemas matemáticos.
    \item \textbf{RdA 2:} Resolver problemas prácticos y teóricos relacionados con el Álgebra Lineal, utilizando técnicas de cálculo con matrices, determinantes y sistemas de ecuaciones, así como en el análisis de espacios vectoriales.
\end{itemize}


%%%%%%%%%%%%%%%%%%%%%%%%%%%%%%%%%%%%%%%%
\subsection*{RdA de la actividad:}
\begin{itemize}[leftmargin=*]
    \item Identificar cuándo una matriz cuadrada es invertible.
    \item Relacionar la existencia de la inversa con el determinante y con la reducción por filas.
    \item Calcular la matriz inversa mediante el método de Gauss-Jordan sobre $(A|I)$.
    \item Verificar resultados analíticos con apoyo de herramientas computacionales.
\end{itemize}

%%%%%%%%%%%%%%%%%%%%%%%%%%%%%%%%%%%%%%%%
\section*{Introducción}
%%%%%%%%%%%%%%%%%%%%%%%%%%%%%%%%%%%%%%%%

%%%%%%%%%%%%%%%%%%%%%%%%%%%%%%%%%%%%%%%%
\paragraph{Control de lectura:}
Antes de iniciar la clase, se aplicará el control de avance:
\href{https://gemini.google.com/share/773ea3fbeaf3}{Control de avance - Clase 02}.

\paragraph{Pregunta inicial:}
En los número reales, la ecuación $ax=b$ tiene solución única $x=a^{-1}b$ siempre que $a\neq 0$. ¿Podemos extender esta idea a matrices? 

%%%%%%%%%%%%%%%%%%%%%%%%%%%%%%%%%%%%%%%%
\section*{Desarrollo}
%%%%%%%%%%%%%%%%%%%%%%%%%%%%%%%%%%%%%%%%

%%%%%%%%%%%%%%%%%%%%%%%%%%%%%%%%%%%%%%%%
\subsection*{Actividad 1: Inversa de una matriz y criterios de invertibilidad}

La clase inicia con una exploración de casos concretos para reconocer el significado de «matriz inversa» y su conexión con el determinante. Luego, mediante clase magistral y práctica guiada, se formaliza el concepto, se estudian propiedades y se aplica el método de Gauss-Jordan para calcular inversas.

\paragraph{¿Cómo lo haremos?}
\begin{itemize}[leftmargin=*]
    \item \textbf{Exploración guiada:} los estudiantes trabajan con dos matrices $2\times 2$ que son inversas entre sí:
    \[
        A = \begin{pmatrix}
            3 & 2 \\
            5 & 3
        \end{pmatrix},\quad
        B = \begin{pmatrix}
            -3 & 2 \\
            5 & -3
        \end{pmatrix}.
    \]
    Se les solicita multiplicarlas en ambos órdenes y justificar por qué en ambos casos se obtiene la identidad. Después, se analiza la matriz
    \[
        \begin{pmatrix}
            1 & 1 \\
            1 & 1
        \end{pmatrix}
    \]
    para discutir, con argumentos, por qué no admite inversa.
    \item \textbf{Clase magistral:} se define formalmente la matriz inversa, sus propiedades algebraicas y su relación con la invertibilidad. Se conecta con la equivalencia por filas de $A$ con $I_n$. Se apoya la explicación en el resumen del curso: \href{https://fcena-puce.github.io/AlgebraLineal-05-N0048/01Resumen/Resumen03.pdf}{Resumen03.pdf}.
    \item \textbf{Resolución de ejercicios:} se desarrollan ejercicios del documento \href{https://fcena-puce.github.io/AlgebraLineal-05-N0048/04Ejercicios/Ejercicios03.pdf}{Ejercicios03.pdf}, enfatizando el cálculo de inversas por el método
    \[
        (A|I_n)\sim(I_n|A^{-1}).
    \]
    \item \textbf{Implementación computacional:} se realiza mediante cuadernos de Jupyter usando el documento \href{https://fcena-puce.github.io/AlgebraLineal-05-N0048/02ResumenPython/ResumenPy03.pdf}{ResumenPy03.pdf}.
\end{itemize}

\paragraph{Verificación de aprendizaje:}
\begin{enumerate}[leftmargin=*]
    \item ¿Qué significa que una matriz $A$ tenga inversa? 
    \item ¿Cómo se relaciona la existencia de la inversa con el determinante de $A$?
    \item ¿Qué método se puede usar para calcular la inversa de una matriz?
\end{enumerate}

%%%%%%%%%%%%%%%%%%%%%%%%%%%%%%%%%%%%%%%%
\section*{Cierre}
%%%%%%%%%%%%%%%%%%%%%%%%%%%%%%%%%%%%%%%%

\paragraph{Tarea:}
Resolver del libro \href{https://catalogobiblioteca.puce.edu.ec/cgi-bin/koha/opac-detail.pl?biblionumber=86081}{Fundamentos de álgebra lineal de Larson}, sección 2.3, los ejercicios: 1, 5, 7, 13, 45, 55.

Además, completar el \href{https://fcena-puce.github.io/AlgebraLineal-05-N0048/05Proyectos-Retos/FactorizacionLU.pdf}{Reto: Factorización $LU$}.

\paragraph{Pregunta de investigación:}
\begin{enumerate}[leftmargin=*]
    \item ¿Qué tan costoso es calcular la inversa de una matriz grande? 
    \item ¿Existen métodos alternativos para resolver sistemas de ecuaciones lineales sin calcular la inversa?
\end{enumerate}

\paragraph{Para la próxima clase:}
Visualizar el video \href{https://www.youtube.com/watch?v=Bd7XboIfsU4}{¿Qué son LOS DETERMINANTES? ¿Son amigos o enemigos?}. Luego realizar la Clase Invertida disponible en \href{https://fcena-puce.github.io/AlgebraLineal-05-N0048/07Act-ClaseInvertida/01ClsInv-Determinantes.pdf}{01ClsInv-Determinantes.pdf}.



\end{document}
